\section{Physical momentum fibres and a complete endpoint kernel}
\label{sec:mc-fibres}
The preceding source, curvature and growth theorems are retained, with
no change to their proof bodies. We use them for a different source:
a fixed initial physical tangential momentum. This is a coordinate
disintegration, not a new conditioning event. In particular we neither
replace $\nu_R^*$ nor remove any collision or section test.
Throughout this section $R\in[R_-,R_+]$, $m\ge1$, and
$a\in(-1,1)$. The source fibre is
\[
       F_a=\{(\alpha_0,p_0):0\le\alpha_0<2\pi,\ p_0=a\},
       \qquad d\lambda_a=d\alpha_0.
\]
Its total mass is $2\pi$, not one. We keep $t_*=3/50$,
$k_-=100/47$, $k_+=20/9$, and $U=170/9$ from
Section~\ref{sec:uc-curvature}. A fixed enlargement of the admissible
geometric-curvature allowance will be explicitly paid below.

\begin{lemma}[First true collision of a momentum fibre]
\label{lem:mc-first}
Put $c=\sqrt{1-a^2}$, $d=\cos\varphi_1$, $p_1=\sin\varphi_1$,
$\kappa=1/R$, $A=c+\tau\kappa$, and $y=d\kappa/A$.
On every actual regular first-hit branch of $F_a$,
\begin{align}
 \frac{ds_1}{ds_0}&=-\frac A d,&
 V_1:=\frac{d\varphi_1}{ds_1}&=\kappa+y,
       \qquad \kappa\le V_1\le U,\label{eq:mc-first}\\
 \frac{dV_1}{ds_1}
   &=-\frac{p_1\kappa}{A}(\kappa+2y)
          -\frac{a d^2\kappa^2}{A^3},\label{eq:mc-second}\\
 \rho_1:=\frac{d\lambda_a}{d\ell_1}
   &=\frac d{R A\sqrt{1+V_1^2}}\le\frac d{t_*}.
                                                        \label{eq:mc-density}
\end{align}
The geometric curvature of this image is bounded by
\begin{equation}\label{eq:mc-curvature}
 K_{\rm mom}:=\frac{k_++2/t_*}{t_*}
                  +\frac1{t_*^3 k_-},
\end{equation}
uniformly in both positive endpoint incidences. This bound is deliberately
coarse; it has no collision-count or microscopic-window dependence.
\end{lemma}
\begin{proof}
The differentiation proving Lemma~\ref{lem:uc-transform} is valid also
for an initial graph of slope zero: its identities precede any use of the
unstable cone. Substitute $V=0$ and $V_s=0$, not $V=\kappa$, in those
identities. This gives \eqref{eq:mc-first}--\eqref{eq:mc-second}.
Alternatively the two projections of $q_1=q_0+\tau v$ give
$-d s_{1,s_0}=c+\tau\kappa$ and
$\tau_{s_0}=-a-p_1A/d$; differentiating the reflection formula gives
the same result. Since $A\ge t_*\kappa$, one has
$y\le1/t_*$ and $\kappa/A\le1/t_*$. Taking absolute values in
\eqref{eq:mc-second}, and then dividing by
$(1+V_1^2)^{3/2}\ge1$, proves \eqref{eq:mc-curvature}.
Finally $d\ell_1=\sqrt{1+V_1^2}|ds_1|$, $ds_0=R\,d\alpha_0$,
and $RA=Rc+\tau\ge t_*$, proving \eqref{eq:mc-density}.
All derivatives are on the actual first-hit component. Nothing is
continued through a physical grazing collision in this statement.
\end{proof}

Fix once and for all
\begin{equation}\label{eq:mc-class}
 \widehat B_{\rm cv}\ge\max\{B_{\rm cv},K_{\rm mom},2C_*\}.
\end{equation}
Theorem~\ref{thm:uc-curvature} and the physical complexity and distortion
proofs apply to this fixed allowance. Let
$\widehat k_0,\widehat\ell_*,\widehat C_g$ be the resulting constants
of Theorem~\ref{thm:uc-growth}. The growth constant is allowed to change;
we do not reuse a constant for a smaller, unverified source class.
We may shorten $\widehat\ell_*$ once. None of these choices depends on
$a$, $m$, the exact record, or the outer parameter $B$.

\begin{lemma}[A complete, grazing-summable first-image family]
\label{lem:mc-family}
For every $a$ the first physical image of $F_a$, with the source labels
retained, decomposes into countably many homogeneous short curves $C$ in
$\mathcal W_R^u(\widehat B_{\rm cv})$, carrying their exact pushed
weights $\rho_C$. It covers $\lambda_a$-almost every initial point and
satisfies
\begin{equation}\label{eq:mc-family}
 \sum_C\|\rho_C\|_\infty\le A_{\rm mom}<\infty
\end{equation}
with a constant uniform in $R$ and $a$.
One explicit choice in terms of the fixed geometric constants is
\begin{align}
 \ell_{\rm pre}
   &=\frac{\pi}{4\{(1-R_+)^{-1}+k_+\}},&
 J_{\rm pre}&=2+\left\lceil\frac{2\pi R_+}{\ell_{\rm pre}}\right\rceil,
                                                        \label{eq:mc-pre}\\
 L_0&=\frac{\pi\sqrt{1+U^2}}{k_-},&
 J_{\rm len}&=2+\left\lceil\frac{2L_0}{\widehat\ell_*}\right\rceil,
                                                        \label{eq:mc-length}\\
 A_{\rm mom}
   &=\frac{J_{\rm pre}J_fJ_{\rm len}}{t_*}
                      \left(1+\frac2{\widehat k_0-1}\right).
                                                        \label{eq:mc-A}
\end{align}
Here $J_f$ and the intrinsic finite $\tau_+$ have exactly the meanings
of \eqref{eq:uc-candidates}; no sampled value is assigned to them.
\end{lemma}
\begin{proof}
Divide the initial boundary circle into at most $J_{\rm pre}$ arcs of
$s_0$-length at most $\ell_{\rm pre}$. For every candidate center $z$
the relative angle $\eta_z=\arg(z-q_0)-\vartheta$ varies by at most
$\pi/4$ on such an arc, and
\[
 \eta_z'\le(1-R_+)^{-1}-k_-<0,\qquad
 b_z'=-c-t_z\kappa.
\]
The positive-foot set is an interval; on it the impact parameter is
strictly decreasing. Thus each disk contributes at most two tangency
cuts. Between these cuts the first-entry times of two disjoint disks
cannot exchange order, since equality would give the same physical
point on both boundaries. The same finite candidate argument as in
Section~\ref{sec:uc-complexity} gives at most $J_f$ first-hit branches
on each preliminary arc. The infinite lattice is reduced here by the
proved finite-horizon bound, not by truncating a computation.

On each branch $V_1\ge k_->0$, so $\varphi_1$ is strictly monotone as a
function of a local lift of $s_1$. Each target homogeneity strip is met
at most once. Since $\varphi_1\in(-\pi/2,\pi/2)$, every such connected
image has length at most $L_0$. Subdivide it into lengths at most
$\widehat\ell_*$, using no more than $J_{\rm len}$ pieces, with the
fixed short-length convention of the growth partition. Circular
coordinate seams are glued, or split once within this conservative
count, not counted as different trajectories.

In the central strip \eqref{eq:mc-density} gives $\rho_C\le t_*^{-1}$;
in each of the two strips of index $k\ge\widehat k_0$ it gives
$\rho_C\le(t_*k^2)^{-1}$. Sum all physical branches, subdivisions and
both strip signs, using
$\sum_{k\ge\widehat k_0}k^{-2}\le(\widehat k_0-1)^{-1}$.
This proves \eqref{eq:mc-A}. The density bound, rather than a finite
number of grazing strips, is essential. Physical tangencies are finite
on each preliminary arc, and strip and length boundaries are countable.
Their preimages have zero $d\alpha_0$ measure. The first hit is unique
everywhere else. Coincident images, if any after a coordinate
identification, retain their different source labels and are summed.
\end{proof}

\begin{lemma}[Exact one-collision lattice kernel and physical roof]
\label{lem:mc-oneflight}
For $a,v\in(-1,1)$ and a candidate lattice displacement
$z\ne0$, put $r_z=|z|$, $\gamma_z=\arg z$ and
\begin{align}
 Q_z(a,v)&=\sqrt{r_z^2-R^2(v-a)^2},\label{eq:mc-Qz}\\
 \alpha_z(a,v)&=\gamma_z-
        \arcsin\!\left(\frac{R(v-a)}{r_z}\right)-\arcsin a
                                                    \pmod{2\pi},\notag\\
 T_z(a,v)&=Q_z(a,v)-R\sqrt{1-a^2}-R\sqrt{1-v^2}.
                                                    \label{eq:mc-one-roof}
\end{align}
Let $\mathcal A_z(a,v)$ be the exact indicator that the regular
physical next collision from $(\alpha_z(a,v),a)$ is with center $z$.
For a bounded Borel source selector $h$, the one-collision endpoint
kernel is, almost everywhere in $v$,
\begin{equation}\label{eq:mc-one-kernel}
 D^h_{1,R}(a,v)=
 \sum_{\substack{z\in\Lambda\\0<|z|\le D_f}}
 h(\alpha_z(a,v),a)\mathcal A_z(a,v)\frac{R}{Q_z(a,v)}.
\end{equation}
The unweighted expression is also the exact sum of regular inverse
root derivatives at each $a,v$. On an admitted term, $T_z$ is the
\emph{original} flight time and $z$ its actual lattice displacement.
In particular
\begin{equation}\label{eq:mc-one-C}
 D_{1,R}(a,v)\le C_1:=\frac{47N_f}{2\sqrt{291}}.
\end{equation}
This is not a Markov transition rule for the momentum process.
\end{lemma}
\begin{proof}
Let $\vartheta=\alpha_0+\arcsin a$ be the outgoing direction.
The signed impact on the next disk is
$b_z=(z-q_0)\cdot v_\vartheta^\perp
     =z\cdot v_\vartheta^\perp+Ra=R p_1$.
Thus $z\cdot v_\vartheta^\perp=R(v-a)$.
The root with $z\cdot v_\vartheta>0$ is exactly the displayed
$\alpha_z$, and $z\cdot v_\vartheta=Q_z$ there. The other angular
root cannot be a physical forward next collision: on a physical
flight this projection is $\tau+R\sqrt{1-a^2}+R\sqrt{1-v^2}>0$.
There is no vanishing angular discriminant, since
$r_z\ge1$ and $2R\le94/100<1$.

At fixed $a$, differentiation gives
$\partial_{\alpha_0}p_1=-z\cdot v_\vartheta/R=-Q_z/R$.
The projection identity just used proves \eqref{eq:mc-one-roof}.
One-dimensional coarea now gives \eqref{eq:mc-one-kernel}, with the
actual next-hit mask retained. The mask excludes a disk hidden by
an earlier collision, not the earlier collision from the physical
record. Every actual first center lies in the proved candidate set,
so the expression is complete. Finally
$Q_z\ge\sqrt{1-4R_+^2}=\sqrt{291}/50$ and
$R/Q_z\le47/(2\sqrt{291})$. Sum at most $N_f$ terms.
The formula keeps the reconstructed initial angle and its exact
physical visibility mask. Discarding that information and composing
these kernels as if the momenta alone were a Markov chain is not
part of this proof or of any subsequent argument.
\end{proof}

\begin{theorem}[All-time endpoint-momentum kernel on the physical source]
\label{thm:mc-kernel}
There exists a finite constant
\begin{equation}\label{eq:mc-C}
 C_{\rm mom}=\max\left\{C_1,
             \frac{A_{\rm mom}\widehat C_gJ_f}{k_-t_*}\right\}
\end{equation}
such that, for every $R$, $m\ge1$, $a\in(-1,1)$, Borel set
$I\subset(-1,1)$, and Borel selector $0\le h\le1$ on the original
physical source,
\begin{equation}\label{eq:mc-kernel}
 \int_{F_a^{\rm reg}(m)}h(\alpha_0,a)
              \one_{\{p_m\in I\}}\,d\alpha_0
       \le C_{\rm mom}|I|.
\end{equation}
The selector may retain the complete original $\Pi_{m,n,k}$, both
section tests, every collision, any original-roof Borel test and any
bounded physical guard. It is not required to factor over time.
For every fixed $a,v\in(-1,1)$ the complete regular root sum obeys
\begin{equation}\label{eq:mc-roots}
 D_{m,R}(a,v):=
 \sum_{\substack{\alpha_0:\ p_m(\alpha_0,a)=v\\
                         (\alpha_0,a)\in F_a^{\rm reg}(m)}}
   \frac1{|\partial_{\alpha_0}p_m(\alpha_0,a)|}
          \le C_{\rm mom}.
\end{equation}
Roots are unique physical initial points, not multiple copies from
homogeneity or coordinate seams. In particular the original section
source has the joint endpoint-momentum domination
\begin{equation}\label{eq:mc-joint}
 (p_0,p_m)_\#(h\,d\nu_R^*)
       \le\frac{C_{\rm mom}}{4\pi c_*}\,da\,dv.
\end{equation}
The statement concerns $\nu_R^*$ as originally defined; no conditional
probability at $p_0=a$ is introduced as a replacement model.
\end{theorem}
\begin{proof}
First expose the cancellation at the \emph{last actual collision}.
For a short incoming unstable graph of slope $V$ and arclength $\ell$,
write $w=\kappa+V$, $A=c+\tau w$ as in
Lemma~\ref{lem:uc-transform}. On each physical next-hit branch,
\begin{equation}\label{eq:mc-last-cancellation}
 \left|\frac{dp_1}{d\ell}\right|
  =\frac{d V_1}{\sqrt{1+V^2}}\left|\frac{ds_1}{ds}\right|
  =\frac{\kappa A+dw}{\sqrt{1+V^2}}
  \ge k_-t_*.
\end{equation}
Indeed $A\ge t_*w$ and
$w/\sqrt{1+V^2}\ge1$ because $\kappa>1$ and $V\ge\kappa$.
The small terminal incidence $d$ cancels exactly before an upper
bound is taken. There are at most $J_f$ physical next-hit branches
on this short graph, each with strictly monotone $p_1$. Do \emph{not}
cut this last coarea into infinitely many terminal homogeneity strips.
Direct one-dimensional substitution therefore gives
\begin{equation}\label{eq:mc-last-kernel}
 \int_V\one_{\{p_1\in I\}}\,d\ell
                      \le\frac{J_f}{k_-t_*}|I|.
\end{equation}

For $m\ge2$, push each first-image density of
Lemma~\ref{lem:mc-family} through $m-2$ additional actual collisions,
using the complete homogeneous short-curve partition. The density on
an incoming final graph is bounded by $\|\rho_C\|_\infty$ times its
exact inverse curve Jacobian. Their complete sum of suprema is at most
$\widehat C_g\|\rho_C\|_\infty$, by
Theorem~\ref{thm:uc-growth}. Now apply \eqref{eq:mc-last-kernel} to
the one remaining physical collision. Sum the labelled graphs and
first-image sources with Tonelli's theorem. This proves the bound
with $A_{\rm mom}\widehat C_gJ_f/(k_-t_*)$.
For $m=1$, directly on each of the at most $J_{\rm pre}J_f$ initial
first-hit pieces,
\[
 \left|\frac{dp_1}{d\alpha_0}\right|
       =R A V_1=A+d\ge t_*\kappa\ge k_-t_*.
\]
Lemma~\ref{lem:mc-oneflight} gives the sharper complete bound $C_1$
used in \eqref{eq:mc-C}. There is neither an
extra initial nor an extra terminal collision in either case.
A selector between zero and one is retained until the positive upper
bound. The first-image and growth singular-null arguments apply on
each labelled source, and the final collision removes only finitely
many tangencies on each incoming short curve. Their pullbacks are null
for $d\alpha_0$ by the exact absolutely continuous weights above.

On each regular original $m$-itinerary the map
$\alpha_0\mapsto p_m(\alpha_0,a)$ has nonzero derivative: after the
first hit it is an unstable graph and all subsequent transformations
are regular. One-dimensional coarea gives \eqref{eq:mc-roots} for
almost every $v$. To obtain the stated bound at every regular root
value, fix a finite collection of distinct roots at $v_0$ and take
disjoint regular source neighborhoods. The inverse branches and their
positive weights are continuous for $v$ near $v_0$, where their sum
is bounded almost everywhere by the already proved density bound.
Pass to $v_0$ and then take the supremum over finite collections.
This proves the bound for the complete nonnegative root series,
including roots on artificial subdivision boundaries. No value of a
physical derivative is assigned at a singular trajectory by this
argument. Finally integrate the fibre inequality in $a$ and use
$d\nu_R^*=\one_{Y_R^*}\,d\alpha_0\,da/(4\pi c_*)$.
The factor $c_*$ occurs exactly once, and there is no $R$ factor
because the initial fibre coordinate is $\alpha_0$, not $s_0$.
\end{proof}

\begin{remark}[What is, and is not, smoothed]
\label{rem:mc-scope}
The bounds sum all relevant regular collision words before any supremum
in $R,m,a,v$ with $a,v\in(-1,1)$, with no endpoint-incidence loss. They have no
whole-cylinder pressure premise. They give no mixing rate, product
law, lattice local limit, or density in the original clock $L_m$.
The grazing-uniform bound follows from the last-collision cancellation
\eqref{eq:mc-last-cancellation}, not from dropping an inverse incidence
factor in a terminal-coordinate coarea estimate. Singular one-sided
critical traces require a separate argument below; they do not follow
by declaring a bulk-null set to have zero trace.
\end{remark}
